\chapter*{Введение}
\addcontentsline{toc}{chapter}{Введение}
\markboth{Введение}{}

Современный этап развития вычислительной техники характеризуется не только последовательным увеличением номинальной производительности процессоров и подсистем памяти, но и заметным усложнением внутренней организации вычислительной платформы. Если ранее оценка возможностей компьютера нередко сводилась к сравнению тактовой частоты, числа ядер или объёма оперативной памяти, то в настоящее время такого описания уже недостаточно. Реальная производительность системы формируется как результат взаимодействия множества компонентов: микроархитектуры центрального процессора, иерархии кэш-памяти, контроллера оперативной памяти, характеристик накопителя, особенностей сетевой подсистемы, политики планировщика операционной системы, фоновой активности служб, а также используемых программных библиотек и алгоритмов. Вследствие этого одна и та же аппаратная конфигурация способна демонстрировать существенно различающееся поведение в зависимости от характера нагрузки, условий эксплуатации и текущего состояния программной среды.

Особенно важным становится тот факт, что вычислительные системы общего назначения всё чаще используются в условиях многопрофильной нагрузки. На одной и той же рабочей станции или сервере могут одновременно выполняться компиляция программ, запуск виртуальных сред, обработка больших объёмов данных, криптографические операции, сетевое взаимодействие и фоновый мониторинг. В такой ситуации возникает практическая потребность не просто в регистрации отдельных метрик, а в их комплексной интерпретации. Пользователю требуется понимать не только «сколько процентов загрузки показывает CPU», но и почему наблюдается именно такая картина, какие компоненты становятся ограничивающими, насколько устойчивы результаты измерений и как связать полученные показатели с реальными прикладными сценариями.

Актуальность темы исследования определяется сразу несколькими обстоятельствами. Во-первых, для разработчиков программного обеспечения и системных инженеров критично иметь инструмент, позволяющий быстро выявлять узкие места вычислительной платформы без необходимости обращаться к большому числу разрозненных утилит. Во-вторых, в образовательной и исследовательской практике возрастает интерес к наглядным средствам демонстрации влияния архитектурных особенностей процессора, памяти и подсистемы ввода-вывода на итоговую производительность. В-третьих, в задачах компьютерной безопасности и защищённой обработки данных всё большее значение имеет оценка криптографической мощности системы, поскольку скорость хеширования, симметричного шифрования и асимметрических операций напрямую влияет на производительность защищённых протоколов, систем аутентификации, средств контроля целостности и сервисов безопасного хранения данных.

Следует учитывать, что традиционные средства операционной системы, такие как диспетчеры задач, системные мониторы и стандартные консольные утилиты, как правило, ориентированы на отображение состояния «в моменте». Они позволяют получить полезные сведения о текущем потреблении ресурсов, однако редко дают целостное представление о вычислительных возможностях платформы. Например, общая загрузка процессора не отражает распределение нагрузки по логическим ядрам и не позволяет судить о характере параллелизма. Метрики памяти часто показывают лишь степень занятости ресурсов, но не дают ответа на вопрос о пропускной способности или латентности доступа. Дисковые и сетевые показатели обычно фиксируются как мгновенные скорости, однако без сопоставления с фоновыми процессами и режимом работы системы их трудно интерпретировать. Аналогично, отдельные бенчмарки позволяют получить итоговый числовой результат, но нередко не объясняют, в каких условиях он получен и как он соотносится с текущим состоянием системы.

Тем самым возникает противоречие между потребностью в комплексном, воспроизводимом и интерпретируемом анализе вычислительной платформы и ограниченностью существующих подходов, которые либо сосредоточены на оперативном мониторинге, либо предоставляют синтетическую оценку без достаточного контекста. Разрешение этого противоречия предполагает создание программного комплекса, объединяющего функции телеметрии и измерительного тестирования в рамках единой архитектуры и общего пользовательского интерфейса. Такой подход позволяет одновременно наблюдать состояние системы, инициировать контролируемую нагрузку и интерпретировать результаты с учётом фоновых факторов, что существенно повышает прикладную ценность анализа.

В данной дипломной работе рассматривается и реализуется программный комплекс «Системный анализатор», предназначенный для исследования вычислительных возможностей микропроцессорных систем общего назначения. Концептуальной основой комплекса является сочетание двух взаимодополняющих режимов работы. Первый режим связан с мониторингом телеметрии в реальном времени и обеспечивает получение информации о ключевых системных характеристиках: параметрах CPU, использовании памяти, активности дисков и сети, средней нагрузке, времени работы системы и перечне наиболее ресурсоёмких процессов. Второй режим представляет собой набор синтетических тестов, сгруппированных по прикладным категориям и предназначенных для оценки предельных или характерных возможностей платформы при вычислительных, память-зависимых, математических и криптографических нагрузках.

Принципиально важным является то, что в рамках предлагаемого подхода мониторинг и бенчмарки не рассматриваются как независимые подсистемы. Напротив, они образуют единую информационную среду анализа. Это означает, что пользователь может сопоставлять результат измерительного теста с текущим фоном системы: активностью других процессов, распределением нагрузки по ядрам, состоянием ввода-вывода и иными факторами, способными искажать измерения. Тем самым повышается интерпретируемость результатов и уменьшается риск ошибочных выводов, возникающих при использовании «изолированных» числовых оценок без контекста.

С научно-практической точки зрения задача анализа вычислительной платформы осложняется многомерностью самого понятия производительности. Для процессорных нагрузок существенны не только количество ядер и частота, но и глубина конвейера, эффективность предсказания ветвлений, работа кэшей, поддержка SIMD-инструкций, механизмы энергосбережения и качество распараллеливания. Для памяти важны одновременно объём, пропускная способность и латентность случайного доступа. Для подсистемы ввода-вывода решающую роль играют характеристики накопителя, тип интерфейса, наличие очередей операций и влияние фоновых служб. Для криптографических задач необходимо учитывать наличие аппаратного ускорения, особенности реализации библиотек и различие в вычислительной природе симметричных, асимметричных и хеш-функциональных алгоритмов. Следовательно, универсальный инструмент анализа должен поддерживать несколько классов измерений и обеспечивать согласованный способ их представления.

Отдельное место в работе занимает проблема воспроизводимости результатов. В многозадачной операционной системе любые синтетические измерения подвержены влиянию случайных возмущений: миграции потоков между ядрами, фоновых обращений к диску, кратковременной сетевой активности, изменения частотных режимов процессора, эффектов прогрева кэшей и иных факторов. Поэтому получение «одного числа» без статистической обработки не может считаться достаточной основой для корректного вывода. В связи с этим в проекте используется подход, предусматривающий разогрев, серию повторных измерений, фильтрацию выбросов и вычисление устойчивых статистических показателей. Такая организация эксперимента позволяет сделать результаты не только более надёжными, но и более полезными в прикладном и исследовательском отношении.

С точки зрения программной реализации важным требованием является кросс-платформенность и практическая применимость решения. Инструмент анализа должен быть пригоден для использования на современных персональных компьютерах и рабочих станциях без сложной настройки, сохранять отзывчивость интерфейса во время выполнения тестов и при этом позволять расширять каталог измерений по мере развития проекта. В данной работе для реализации выбран язык программирования Go, сочетающий нативную компиляцию, развитую стандартную библиотеку, удобные средства конкурентного выполнения и хорошие возможности для системного программирования. Для построения графического интерфейса применяется библиотека Fyne, обеспечивающая единый API для разных операционных систем и достаточную гибкость при создании вкладок, панелей и интерактивных элементов управления.

Архитектурно программный комплекс разделён на несколько логически самостоятельных модулей. Подсистема \texttt{pkg/profiling} отвечает за сбор и агрегирование телеметрии операционной системы, включая сведения о процессоре, памяти, дисках, сети, времени работы, средней нагрузке и активных процессах. Подсистема \texttt{pkg/benchmark} реализует каталог измерительных тестов, их унифицированный интерфейс, механизм запуска и статистическую обработку результатов. Подсистема \texttt{pkg/gui} обеспечивает пользовательское представление данных, организацию вкладок, визуализацию состояния системы и сценарии запуска тестов, в том числе интерактивные криптографические измерения с изменяемыми параметрами. Подобное разбиение не только повышает сопровождаемость кода, но и соответствует методологической задаче отделения получения данных, вычислительной логики и пользовательского представления.

Следует подчеркнуть, что объектом внимания в работе является не абстрактный «идеальный» бенчмарк, а инструмент, ориентированный на практику системного анализа. Это означает, что при проектировании учитываются не только требования к точности измерений, но и требования к интерпретируемости, удобству использования и связи результатов с прикладными сценариями. Пользователю важно не просто узнать, что та или иная система показывает определённое число мегабайт в секунду или операций в секунду, а понять, что именно означает данный показатель, при каких предпосылках он получен, насколько он стабилен и какие реальные выводы можно из него сделать относительно пригодности платформы для разработки, анализа данных, криптографической обработки или эксплуатации ресурсоёмких сервисов.

\section*{Цель, задачи и практическая значимость}

\textbf{Целью дипломной работы} является разработка, теоретическое обоснование и практическая реализация программного комплекса для комплексного анализа вычислительных возможностей микропроцессорных систем, обеспечивающего одновременное получение телеметрии в реальном времени и воспроизводимых результатов синтетического тестирования по нескольким категориям нагрузок.

Для достижения поставленной цели в работе решаются следующие \textbf{задачи}:
\begin{enumerate}
    \item Выполнить анализ предметной области и обосновать необходимость объединения средств мониторинга и бенчмаркинга в рамках единого инструмента системного анализа.
    \item Исследовать ключевые классы метрик, характеризующих вычислительную платформу, включая процессорные, память-зависимые, дисковые, сетевые и криптографические показатели.
    \item Сформулировать функциональные и нефункциональные требования к программному комплексу с учётом сценариев практической эксплуатации: первичной диагностики, сравнения платформ, учебного применения и анализа аппаратных ограничений.
    \item Спроектировать модульную архитектуру приложения и выбрать технологический стек, обеспечивающий кросс-платформенность, расширяемость, достаточную производительность измерений и отзывчивость пользовательского интерфейса.
    \item Реализовать подсистему телеметрии, включающую сбор сведений о процессоре, памяти, дисковой и сетевой активности, времени работы системы, средней нагрузке и списке активных процессов.
    \item Реализовать библиотеку синтетических тестов по категориям «Процессор», «Память», «Математика», «Криптография» и обеспечить единый механизм их запуска и интерпретации.
    \item Разработать подход к повышению устойчивости результатов, включающий разогрев, многократные прогоны, фильтрацию выбросов и вычисление робастных статистических показателей.
    \item Создать графический интерфейс, позволяющий пользователю наблюдать состояние системы, запускать тесты, анализировать результаты и выполнять интерактивные криптографические измерения с изменяемыми параметрами.
    \item Провести анализ применимости разработанного комплекса, описать ограничения методики измерений и сформулировать рекомендации по интерпретации получаемых результатов.
\end{enumerate}

\textbf{Объект исследования} --- вычислительные системы общего назначения, включающие многоядерные микропроцессоры, иерархию памяти, типовые подсистемы ввода-вывода и программную среду, определяющую условия выполнения прикладных задач.

\textbf{Предмет исследования} --- методы программного мониторинга и синтетического тестирования вычислительной платформы, способы повышения воспроизводимости результатов, а также подходы к представлению и интерпретации измеряемых показателей в задачах системного анализа и компьютерной безопасности.

\textbf{Практическая значимость} работы заключается в создании прикладного программного инструмента, который может быть использован:
\begin{itemize}
    \item для первичной диагностики производительности и текущего состояния рабочей станции, ноутбука или сервера перед развёртыванием ресурсоёмких приложений и сервисов;
    \item для сравнительного анализа аппаратных платформ по различным профилям нагрузок, включая вычисления общего назначения, операции с памятью, ввод-вывод и криптографическую обработку;
    \item для выявления факторов нестабильности измерений, обусловленных фоновой активностью, конкуренцией за ресурсы и особенностями планировщика операционной системы;
    \item в учебном процессе как наглядное средство демонстрации взаимосвязи между архитектурой вычислительной системы, состоянием программной среды и наблюдаемыми характеристиками производительности;
    \item в задачах предварительной оценки пригодности вычислительной платформы для приложений, чувствительных к пропускной способности памяти, дисковой активности или криптографической нагрузке.
\end{itemize}

\textbf{Научная новизна и методическая ценность} работы в прикладном контексте заключаются в том, что анализ вычислительной платформы рассматривается не как сумма разрозненных метрик, а как единый процесс сопоставления телеметрии и результатов контролируемых измерений. Подобный подход позволяет перейти от простого наблюдения к осмысленной интерпретации, что особенно важно в условиях многозадачных систем, где числовой результат без контекста может быть недостаточно информативным или даже вводящим в заблуждение.

\textbf{Структура работы} отражает логику исследования. Во введении обосновываются актуальность темы, цель, задачи, объект и предмет исследования, а также практическая значимость работы. В первой главе рассматриваются теоретические основы системного анализа вычислительных платформ и анализ существующих подходов. Во второй главе формулируются требования к комплексу и описываются архитектурные решения, положенные в основу программной реализации. В третьей главе анализируются особенности реализации модулей мониторинга, библиотеки тестов, статистической обработки результатов и пользовательского интерфейса. В четвёртой главе описывается методика проведения экспериментов и подход к анализу результатов. В заключении подводятся итоги выполненной работы и намечаются направления дальнейшего развития проекта.

Таким образом, разработка программного комплекса «Системный анализатор» представляет собой актуальную и практически значимую задачу, находящуюся на пересечении системного программирования, анализа производительности, прикладного бенчмаркинга и компьютерной безопасности. Решение данной задачи позволяет получить инструмент, который не только фиксирует параметры работы системы, но и помогает интерпретировать их в контексте реальных вычислительных сценариев, что делает результаты анализа более содержательными, воспроизводимыми и полезными для пользователя.
